%------------------------------------------------------------------------------%
\section{Sistemas de Coordenadas}
\label{sec:sistemas_coorenadas}

O sistema de \emph{Coordenadas Polares}\index{Coordenadas polares}
é definido a partir do sistema de coordenadas cartesiano $(x, y)$
pela transformação de coordenadas
\begin{equation}
  \label{eq:coordenadas_polares}
  \begin{split}
    x & = r \cos(\theta) \\
    y & = r \sin(\theta)
  \end{split}
\end{equation}
onde um ponto $P$ com coordenadas $(x, y)$ pode ser representado também
pelas coordenadas $(r, \theta)$, como mostra a Figura~\ref{fig:coordenadas_polares}.
A transformação inversa é dada por
\begin{equation}
  \label{eq:coordenadas_polares_inversa}
  \begin{split}
    \theta & = \arctan\left(\dfrac{\,y\,}{x}\right) \\[3mm]
    r      & = \sqrt{x^2 + y^2\,}
  \end{split}
\end{equation}
Note que todo o plano cartesiano $(x, y)\in\R^2$ é mapeado na região
\[
  \begin{split}
    r      & \in [\,0, \infty) \\
    \theta & \in [\,0, 2\pi)
  \end{split}
\]
onde o intervalo para $\theta$ pode ser substituído por qualquer outro
intervalo de comprimento $2\pi$.
Apesar da origem $(0, 0)$ ser um ponto ordinário no plano, sua representação
em coordenadas polares é singular, pois se $r=0$ todos os valores de $\theta$
representam o mesmo ponto.

%------------------------------------------------------------------------------%
\begin{figure}
  \centering
  % Title: gl2ps_renderer figure
% Creator: GL2PS 1.4.2, (C) 1999-2020 C. Geuzaine
% For: Octave
% CreationDate: Fri Sep 10 20:09:40 2021
\setlength{\unitlength}{1pt}
\begin{picture}(0,0)
\includegraphics[width=10cm,height=6cm]{fig_coordenadas_polares}
\end{picture}%
\begin{picture}(283,170)(0,0)
\fontsize{12}{0}\selectfont\put(220.697,18.144){\makebox(0,0)[t]{\textcolor[rgb]{0,0,0}{{$x$}}}}
\fontsize{12}{0}\selectfont\put(44.5444,139.629){\makebox(0,0)[r]{\textcolor[rgb]{0,0,0}{{$y$}}}}
\fontsize{12}{0}\selectfont\put(111.361,48.5152){\makebox(0,0)[l]{\textcolor[rgb]{0,0,0}{{$\theta$}}}}
\fontsize{12}{0}\selectfont\put(135.658,106.22){\makebox(0,0){\textcolor[rgb]{0,0,0}{{$r$}}}}
\fontsize{12}{0}\selectfont\put(226.772,145.703){\makebox(0,0)[l]{\textcolor[rgb]{0,0,0}{{$P$}}}}
\end{picture}

  \caption{Sistema de coordenadas polares}
  \label{fig:coordenadas_polares}
\end{figure}
%------------------------------------------------------------------------------%

As derivadas de uma função em coordenadas cartesianas e polares
estão ligadas pela regra da cadeia. Considere a função
\[
  g(x, y) = f\big( r(x, y), \theta(x, y) \big)
\]
pela regra da cadeia em duas dimensões sabemos que
\begin{align*}
  \frac{\partial g}{\partial x}
  & = \frac{\partial f}{\partial r     } \frac{\partial r     }{\partial x}
    + \frac{\partial f}{\partial \theta} \frac{\partial \theta}{\partial x} \\[5mm]
  \frac{\partial g}{\partial y}
  & = \frac{\partial f}{\partial r     } \frac{\partial r     }{\partial y}
    + \frac{\partial f}{\partial \theta} \frac{\partial \theta}{\partial y}
\end{align*}
Calculando as derivadas da transformação~\eqref{eq:coordenadas_polares_inversa} temos
\begin{equation}
\label{eq:coordenadas_polares_regra_cadeia}
\begin{split}
  \frac{\partial g}{\partial x}
  & = \cos\theta           \,\frac{\partial f}{\partial r     }
    - \frac{\sin\theta}{r} \,\frac{\partial f}{\partial \theta}
     \\[5mm]
  \frac{\partial g}{\partial y}
  & =  \sin\theta           \,\frac{\partial f}{\partial r     }
    +  \frac{\cos\theta}{r} \,\frac{\partial f}{\partial \theta}
\end{split}
\end{equation}
Note que nessas expressões $r$ e $\theta$ são pensados como funções de
$x$ e $y$, isso é, $r=r(x, y)$ e $\theta = \theta(x, y)$.
Podemos reproduzir esses cálculos para obter relação recíproca
\begin{equation}
  \label{eq:coordenadas_polares_regra_cadeia_reciproca}
  \begin{split}
    \frac{\partial f}{\partial r}
    & = \cos\theta \,\frac{\partial g}{\partial x}
      + \sin\theta \,\frac{\partial g}{\partial y}
    \\[5mm]
    \frac{\partial f}{\partial \theta}
    & =
      - r \sin\theta \,\frac{\partial g}{\partial x}
      + r \cos\theta \,\frac{\partial g}{\partial y}
  \end{split}
\end{equation}
Uma aplicação das coordenadas polares é a integração em regiões que são mais
facilmente descritas nessas coordenadas. Por exemplo, o disco de raio $\rho$
centrado na origem, $D$, é descrito em coordenadas cartesianas como
\[
  D = \left\{
        (x, y) \in\R^2 \st
        -\rho \leq x \leq \rho
        \quad\text{ e }\quad
        -\sqrt{ \rho^2 -x^2} \leq y \leq \sqrt{ \rho^2 - x^2}\,
  \right\}
\]
enquanto que em coordenadas polares ele pode ser descrito como
\[
  D = \left\{
  (r, \theta) \in\R^2 \st
    0 \leq x \leq \rho
    \quad\text{ e }\quad
    0 \leq \theta < 2\pi \,
  \right\}
\]
Assim a integral de uma função $f$ nessa região assume a forma
\[
  I = \int_D f ds
    = \int_{-\rho}^{\rho} \int_{-\sqrt{ \rho^2 -x^2}}^{\sqrt{ \rho^2 -x^2}} \, f(x, y) \, dy dx
    = \int_0^{2\pi} \int_0^\rho  f(r, \theta) \,r\, dr d\theta
\]
Cuidado para não esquecer que a mudança de variáveis transforma o diferencial
$dxdy$ em $rdrd\theta$.


%------------------------------------------------------------------------------%
